\documentclass[10pt,a4paper]{amsart}
\usepackage{amsmath,amssymb,amsthm,graphicx}
\usepackage{amsmath}
\usepackage{amssymb}
\usepackage{hyperref}
\renewcommand{\theenumi}{\Alph{enumi}}
\renewcommand{\thesubsection}{\thesection\Alph{subsection}}

%%%%%%%%%%%%%%%%%%%%%%%%%%%%%%%%%%%%%%%%%%
% ENVIRONMENTS
%%%%%%%%%%%%%%%%%%%%%%%%%%%%%%%%%%%%%%%%%%%
\newtheorem{theorem}{Theorem}[section]
\newtheorem{prop}[theorem]{Proposition}
\newtheorem{lemma}[theorem]{Lemma}
\newtheorem{remark}[theorem]{Remark}
\newtheorem{corollary}[theorem]{Corollary}
\newtheorem{definition}[theorem]{Definition}
\newtheorem{example}[theorem]{Example}
%
\def\beq{\begin{equation}}
\def\eeq{\end{equation}}
\def\be{\begin{equation*}}
\def\ee{\end{equation*}}
%%%%%%%%%%%%%%%%%%%%%%%%%
% Beginning of the text %
%%%%%%%%%%%%%%%%%%%%%%%%%
%\begin{document}



%\documentclass[10pt]{article}
\usepackage{amssymb}
\usepackage{amsfonts}
%\topmargin 0.3in
%\textwidth 6in
%\textheight 9in%\baselineskip
%\hoffset=-50pt
%\voffset=-2pt
%\newenvironment{proof}[1][Proof]
%{\ \rule{0.5em}{0.5em}}


\title{Dedekind-finite Cardinals having countable partitions}
\author{Supakun Panasawatwong and J K Truss} 
\date{This paper is an amplified version of part of the first author's PhD thesis at the University of Leeds}%\footnotemark}
\title{The Dedekind-finite tree U\_T falls into the definable cardinality class Delta\_5}
\author{Supakun Panasawatwong, J. K. Truss}
\date{arxiv:2204.05700v1; CC BY 4.0}
\begin{document}
\maketitle
\noindent\textit{Editorial note.} All delivered mathematical text is the author's own, copied from the pinned publisher source sources/dedekind-finite-trees-2204.05700v1/source0.tex; no statement or proof step was rewritten, reordered or removed.\par\medskip
\noindent\textit{Editorial note.} Preamble provenance: the retained preamble is the source's own layout adapter, if any, plus the author's macro block at lines 5-44 of sources/dedekind-finite-trees-2204.05700v1/source0.tex (everything between the publisher class metadata and \textbackslash{}begin\{document\}, comments excluded). Nothing was added, no macro was invented, and the publisher document class is deliberately not reproduced; the excerpt serves only as preamble, never as a span coordinate.\par\medskip
\noindent\textit{Editorial note.} External references are navigation pointers only; every deliberately kept step of the author's argument stays inside the counted spans. The displayed numbers of those pointers are the author's own numbering under the author's \textbackslash{}newtheorem declarations (the source has no subsections, so numbering is section.index), computed mechanically by the lane's scanner.\par\medskip
\noindent\textit{Editorial note.} The counted proof contains no reference to any labelled item and no bibliography citation, so the card needs neither a navigation stub nor a bibliography excerpt.\par\medskip
\noindent\textit{Editorial note.} This is the first card of this lane with route logic, following the Chief's ruling that the route field records the content of the source rather than a restriction of the lane.\par\medskip
\noindent\textit{Editorial note.} Source metadata: this source directory stages no SOURCE.json, so the title and the author list were taken from the publisher abstract page kept in the same directory (sources/dedekind-finite-trees-2204.05700v1/arxiv-abs.html, the citation\_title and citation\_author meta entries), and the identifiers from the staged licence evidence. r001 was rejected by the assembler with 'complete author list required' before any compile.\par\medskip
\noindent\textit{Editorial note.} Licence binding: the source stages its licence as LICENSE\_EVIDENCE.json, which carries no status field, so the card points at a per-card reuse record in its own task directory (license-reuse.json) that names the licence, its version URL, the retained source notice it derives from and the fact that no new licence claim is made. r002 pointed at the source's file directly and was reported by the read-only precheck as 'no supported status and no retained source notice'.\par\medskip
\noindent\textit{Allowed standard prerequisites (editorial).}\begin{itemize}
\item The FM-model with finite supports built from the automorphism group of the tree, and the definable closure operation on it. Set up in the paragraphs before the counted statement (source lines 303-305) and used throughout the counted proof. Supplies the model in which the cardinality comparison is made and the closure operator with which the support of a function is measured.
\item The tree-theoretic description of [A union \{x\}] as [A] together with the new point and the meets it creates. Derived inside the counted proof at lines 317-319. Is what lets the two case analyses conclude that a point outside [A] is fixed by the function.
\end{itemize}\medskip
\begin{theorem}  \label{3.1} For any countable weakly $2$-transitive tree $T$, $|U_T| \in \Delta_5$ in ${\mathfrak N}_T$, 
and if $T$ has infinitely many ramification orders, then $|U_T| \not \in \Delta_4$.   \end{theorem}

\begin{proof} Let $f: U_T \to U_T \cup \{*\}$ where $* \not \in U_T$ be surjective in ${\mathfrak N}_T$ with finite 
support $A$.

We remark that for any $x, y$, if $f(x) = y$, then $y \in [A \cup \{x\}]$. For since $A$ supports $f$, $\pi(f) = f$ 
for all $\pi \in G_A$. Then $\langle x, y \rangle \in f$ and so $\langle \pi x, \pi y \rangle \in \pi(f) = f$. Hence if 
$\pi y \neq y$, then $\pi x \neq x$ since $f$ is a function. If $y \not \in [A \cup \{x\}]$, then by definition of 
definable closure $[ \;\;]$, there is $\pi \in G_{A \cup \{x\}} \subseteq G_A$ such that $\pi y \neq y$ but $\pi x = x$, 
contrary to what we have shown. Hence $y \in [A \cup \{x\}]$. Also note that $[A \cup \{x\}]$ is the union of 
$[A] \cup \{x\}$ and $\{x \wedge a: a \in A\}$, together with the minimal points of cones at exceptional ramification 
points in this set. 

Now consider $y \not \in [A]$ and let $x$ be such that $f(x) = y$. We aim to show that $x = y$. Now $x \in [A]$ would imply 
that $y \in [A]$, contrary to supposition, and we deduce that $x \not \in [A]$.

\noindent{\bf Case 1}: $y \not \le a$ for all $a \in [A]$.

Then $y \neq x \wedge a$ for all $a \in A$, and also $y$ is not the minimal point of a cone at an exceptional ramification 
point of $[A]$ (as this would lie in $[A]$). So if $y \neq x$, then $y$ must be the minimal point of a cone at an exceptional
ramification point of the form $x \wedge a$ for some $a \in A$. 

Let $z$ be such that $f(z) = x$. If $x \le b$ for some $b \in [A]$, then $x \wedge a = b \wedge a$, which is impossible.
Applying the same argument as above for $x$ and $z$, $x$ must be the minimal point of a cone at an exceptional ramification
point of the form $z \wedge b$ for some $b \in A$. But now again $x \wedge a = b \wedge a$ which is impossible.

The conclusion is that $x = y$.

\noindent{\bf Case 2}: $y \le a$ for some $a \in [A]$.

If $x \not \le b$ for all $b \in [A]$, then by the argument in Case 1, $f(x) = x$, and $x = y$ again follows. Hence $x \le b$
for some $b \in A$. Then for any $c \in A$, either $x \wedge c = x$ or $x \wedge c \in [A]$. Since $x \not \in [A]$, no new
ramification points are added in passing from $[A]$ to $[A \cup \{x\}]$, and therefore $[A \cup \{x\}] = [A] \cup \{x\}$. It 
follows that $x = y$. 

From both cases, we have $f(y) = y$ for all $y \not \in [A]$. Hence $f[A] = [A] \cup \{*\}$ which is a 
contradiction since $f$ is surjective but $[A]$ is finite. Therefore such $f$ does not exist
in ${\mathfrak N}_T$ and so $|U_T| \in \Delta_5$ in ${\mathfrak N}_T$.  

The final clause follows from the fact that $T$ is not $\aleph_0$-categorical.  \end{proof}

\end{document}
